% Source: https://artofproblemsolving.com/community/c3h1980055_a_question_about_relating_probability_with_geometry
% Stable identity: aops-topic-1980055-q01
A rectangle has side lengths $1$ and $2$. On each of the two sides of length $2$, a semicircle of radius $1$ is drawn inside the rectangle, centred at the midpoint of that side. A point is chosen uniformly at random in the rectangle.

Two students compute the probability that the point lies inside both semicircles. Student A argues: the probability of landing in one semicircle is $\dfrac{\pi}{2}\div2=\dfrac{\pi}{4}$, so the probability of landing in both is $\left(\dfrac{\pi}{4}\right)^2=\dfrac{\pi^2}{16}\approx0.617$. Student B computes the area of the overlap and gets $\dfrac{\pi}{3}-\dfrac{\sqrt3}{4}\approx0.614$. Which answer is correct, and what is wrong with the other argument?
